\chapter{Shortest paths}\label{ch:sp}
\index{shortest path}\index{Dijkstra}\index{Bellman--Ford}\index{Floyd--Warshall}

Four algorithms cover every shortest-path problem you will meet:
BFS when edges are unweighted, Dijkstra when weights are
non-negative, Bellman--Ford when weights can go negative, and
Floyd--Warshall when you need every pair at once. Each is built from
one operation --- relaxation --- and each is right exactly when its
invariant holds; Dijkstra even fails in a specific, drawable way when
the invariant breaks. After this chapter you will pick the algorithm
by the weights and the question, write all four from memory, detect
and reconstruct negative cycles, and fold conditions like ``one
coupon'' into the graph itself.

\section{Relaxation}\label{sec:relax}\index{relaxation}
The atomic move: $d[v] \leftarrow \min(d[v],\, d[u] + w(u,v))$. Every
algorithm here is a different strategy for \emph{ordering} and
\emph{counting} relaxations.
\begin{figure}[ht]\centering
\begin{tikzpicture}[font=\footnotesize]
\node[circle,draw=slate,font=\tiny] (u) at (0,0) {u};
\node[circle,draw=slate,font=\tiny] (v) at (2.4,0) {v};
\draw[-{stealth},slate] (u)--(v) node[midway,above,font=\tiny] {w=4};
\node[font=\tiny] at (0,0.65) {dist=7};
\node[font=\tiny,anchor=west] at (1.7,0.65) {dist=15 $\to$ 11};
\node[font=\tiny] at (1.2,-0.7) {$7+4 < 15$: relax};
\end{tikzpicture}

\caption{One relaxation step. Notice the check is strict inequality
--- an equal-length alternative changes nothing.}
\label{fig:relax}
\end{figure}

\section{BFS, revisited}\label{sec:bfs2}
With every weight 1, the BFS layers of Chapter~\ref{ch:graphs}
\emph{are} shortest paths: the queue is a degenerate priority queue.
Keep BFS for unweighted problems --- it is $O(V{+}E)$ and Dijkstra on
unit weights buys nothing.

\section{Bellman--Ford}\label{sec:bf}\index{Bellman--Ford}
Relax \emph{every} edge, $n{-}1$ times. Correctness: after round $k$
every shortest walk of at most $k$ edges is final, and any shortest
path has $\le n{-}1$ edges. A further improvement in round $n$
certifies a negative cycle --- reachable or not, if you seed every
distance with 0 (an implicit super-source).
\begin{lstlisting}[language=C++,caption={\texttt{bellmanford.h}:
rounds, parent pointers, one-cycle rebuild.},label={lst:bf}]
\end{lstlisting}
\lstinputlisting{code/bellmanford.h}
\begin{example}[CSES \emph{Cycle Finding}]\label{ex:cycfind}
Report any negative cycle.
\par\smallskip
\textbf{Strategy.} Run the rounds; a vertex still improving in round
$n$ is downstream of a cycle. Walking parents $n$ steps lands on the
cycle itself; collect until repeat.
\par\smallskip
\begin{solution}
\lstinputlisting{code/cycle_finding.cpp}
Real output (one valid cycle):
\begin{lstlisting}[numbers=none]
YES
4 1 2 4
\end{lstlisting}
\end{solution}
\end{example}

\section{Dijkstra}\label{sec:dij}\index{Dijkstra}
Greedy frontier: pop the unsettled vertex with smallest tentative
distance; that value is final. Correctness is a cut argument ---
settled vertices sit on one side, and the cheapest unsettled vertex
cannot improve through any path that leaves the settled set, because
all such detours cost at least as much (weights $\ge 0$ is the
load-bearing clause). Drop that clause and the argument dies:
\begin{figure}[ht]\centering
% fig37-negfail.tex -- Figure 7.9: where Dijkstra dies: s->a 2, s->b 5, b->a -4.
% fstyle.tex -- shared TikZ styles for the figures of this book (palette from lifeos.sty).
\tikzset{%
  cell/.style={draw=ink, minimum width=7mm, minimum height=6.5mm, inner sep=0pt, font=\footnotesize, fill=white},
  cellh/.style={cell, fill=accentsoft, draw=accent},
  cella/.style={cell, fill=amberbg, draw=amber},
  cellg/.style={cell, fill=goodbg, draw=good},
  cellx/.style={cell, fill=graybg, draw=slate, text=slate},
  vtx/.style={circle, draw=ink, minimum size=6.5mm, inner sep=0pt, font=\footnotesize, fill=white},
  vtxh/.style={vtx, draw=accent, fill=accentsoft},
  vtxa/.style={vtx, draw=amber, fill=amberbg},
  vtxg/.style={vtx, draw=good, fill=goodbg},
  vtxd/.style={vtx, draw=slate, fill=graybg, text=slate},
  lab/.style={font=\footnotesize, text=slate},
  labA/.style={font=\footnotesize, text=accent},
  labM/.style={font=\footnotesize, text=amber},
  idx/.style={font=\tiny, text=slate},
  arr/.style={-{stealth}, thick, ink},
  arrA/.style={-{stealth}, thick, accent},
  arrM/.style={-{stealth}, thick, amber},
  arrG/.style={-{stealth}, thick, good},
  edge/.style={thick, ink},
  edgeA/.style={thick, accent},
  edgeM/.style={thick, amber},
}%
% \cellrow{prefix}{x0}{y}{style}{v0,v1,...}: a row of cells named prefix0, prefix1, ... (pitch 0.75cm)
\providecommand{\cellrow}[5]{\foreach \v [count=\i from 0] in {#5} {\node[#4] (#1\i) at (#2+\i*0.75,#3) {\v};}}%
% \cellidx{prefix}{n}{start}: index labels under cells prefix0 .. prefix(n-1)
\providecommand{\cellidx}[3]{\foreach \i in {0,...,\numexpr#2-1\relax} {\node[idx, below=0pt of #1\i] {\the\numexpr\i+#3\relax};}}%
% \fitfig{<tikzpicture>}: shrink the picture to \linewidth only when it would overflow.
\ifdefined\fitbox\else\newsavebox{\fitbox}\fi%
\providecommand{\fitfig}[1]{\sbox{\fitbox}{#1}\ifdim\wd\fitbox>\linewidth\resizebox{\linewidth}{!}{\usebox{\fitbox}}\else\usebox{\fitbox}\fi}%

\begin{tikzpicture}[font=\footnotesize]
\node[vtx, minimum size=8mm] (s) at (0,0) {$s$};
\node[vtxd, minimum size=8mm] (a) at (3,1.3) {$a$};
\node[vtx, minimum size=8mm] (b) at (3,-1.3) {$b$};
\draw[arr] (s) -- (a) node[midway, above left, lab] {2};
\draw[arr] (s) -- (b) node[midway, below left, lab] {5};
\draw[arrM] (b) -- (a) node[midway, right, labM] {$-4$};
\node[lab, anchor=west, align=left] at (4.3,1.3) {Dijkstra settles $a$ at $2$ and never looks back};
\node[labM, anchor=west, align=left] at (4.3,-1.3) {truth: $s\to b\to a$ costs $5-4=1$};
\end{tikzpicture}

\caption{Where Dijkstra dies, drawn. Notice $s$ settles $a$ at 2
before $b$ is ever examined --- no negative-edge signal can pull it
back.}
\label{fig:negfail}
\end{figure}
\begin{figure}[ht]\centering
\begin{tikzpicture}[font=\footnotesize]
\begin{scope}\node[circle,draw=slate,fill=accentsoft,font=\tiny] at (0,0) {s};
\node[circle,draw=slate,font=\tiny] at (1,0.6) {}; \node[circle,draw=slate,font=\tiny] at (1,-0.6) {};
\node[font=\tiny] at (0.5,-1.2) {settle s};\end{scope}
\begin{scope}[xshift=2.2cm]\node[circle,draw=slate,fill=graybg,font=\tiny] at (0,0) {s};
\node[circle,draw=slate,fill=accentsoft,font=\tiny] at (1,0.6) {}; \node[circle,draw=slate,font=\tiny] at (1,-0.6) {};
\node[font=\tiny] at (0.5,-1.2) {settle near};\end{scope}
\begin{scope}[xshift=4.4cm]\node[circle,draw=slate,fill=graybg,font=\tiny] at (0,0) {s};
\node[circle,draw=slate,fill=graybg,font=\tiny] at (1,0.6) {}; \node[circle,draw=slate,fill=accentsoft,font=\tiny] at (1,-0.6) {};
\node[font=\tiny] at (0.5,-1.2) {settle next};\end{scope}
\begin{scope}[xshift=6.6cm]\node[circle,draw=slate,fill=graybg,font=\tiny] at (0,0) {s};
\node[circle,draw=slate,fill=graybg,font=\tiny] at (1,0.6) {}; \node[circle,draw=slate,fill=graybg,font=\tiny] at (1,-0.6) {};
\node[font=\tiny] at (0.5,-1.2) {done};\end{scope}
\end{tikzpicture}

\caption{The settled set growing. Notice each frame fixes one vertex
permanently --- that permanence is both the speed and the
vulnerability.}
\label{fig:dijset}
\end{figure}
\begin{lstlisting}[language=C++,caption={\texttt{dijkstra.h}: lazy
deletion via the stale-entry skip.},label={lst:dij}]
\end{lstlisting}
\lstinputlisting{code/dijkstra.h}
\begin{example}[CSES \emph{Shortest Routes I}]\label{ex:routes1}
Distances from city 1.
\par\smallskip
\textbf{Strategy.} Non-negative weights name Dijkstra; skip stale
queue entries instead of deleting them.
\par\smallskip
\begin{solution}
\lstinputlisting{code/shortest_routes1.cpp}
Real output (CSES sample):
\begin{lstlisting}[numbers=none]
0 5 2
\end{lstlisting}
\end{solution}
\end{example}

\section{Floyd--Warshall}\label{sec:fw}\index{Floyd--Warshall}
All pairs at once: $d[i][j] \leftarrow \min(d[i][j],\, d[i][k]{+}d[k][j])$
with $k$ the \emph{outermost} loop --- the invariant is ``only
vertices $1..k$ may appear internally.'' $O(n^3)$ is right when
$n \le 500$.
\begin{lstlisting}[language=C++,caption={\texttt{floyd.h}: the
15-line core + \texttt{next} table.},label={lst:fw}]
\end{lstlisting}
\lstinputlisting{code/floyd.h}
\begin{example}[CSES \emph{Shortest Routes II}]\label{ex:routes2}
$q$ distance queries.
\par\smallskip
\textbf{Strategy.} Queries make it all-pairs; $n\le500$ makes the
cube free.
\par\smallskip
\begin{solution}
\lstinputlisting{code/routes2.cpp}
Real output (CSES sample):
\begin{lstlisting}[numbers=none]
5
5
8
-1
3
\end{lstlisting}
\end{solution}
\end{example}
\begin{figure}[ht]\centering
\begin{tikzpicture}[font=\tiny]
\begin{scope}
  \node[anchor=west] at (0,0) {0\quad 4\quad $\infty$\quad 2};
  \node[anchor=west] at (0,-0.35) {5\quad 0\quad 7\quad 3};
  \node[anchor=west] at (0,-0.7) {$\infty$\quad $\infty$\quad 0\quad 1};
  \node[anchor=west] at (0,-1.05) {6\quad 1\quad 8\quad 0};
  \node at (0.8,-1.7) {before};
\end{scope}
\begin{scope}[xshift=2.9cm]
  \node[anchor=west] at (0,0) {0\quad 4\quad $\infty$\quad 2};
  \node[anchor=west] at (0,-0.35) {5\quad 0\quad 7\quad 3};
  \node[anchor=west] at (0,-0.7) {$\infty$\quad $\infty$\quad 0\quad 1};
  \node[anchor=west] at (0,-1.05) {6\quad 1\quad 8\quad 0};
  \node at (0.8,-1.7) {after $k=1$};
\end{scope}
\begin{scope}[xshift=5.8cm]
  \node[anchor=west] at (0,0) {0\quad 4\quad 11\quad 2};
  \node[anchor=west] at (0,-0.35) {5\quad 0\quad 7\quad 3};
  \node[anchor=west] at (0,-0.7) {$\infty$\quad $\infty$\quad 0\quad 1};
  \node[anchor=west] at (0,-1.05) {6\quad 1\quad 8\quad 0};
  \node at (0.8,-1.7) {after $k=2$};
\end{scope}
\node[anchor=west] at (0,-2.4) {each stage lets routes detour through
one more vertex};
\end{tikzpicture}

\caption{The $k$-expansion. Notice each stage only \emph{permits}
detours through one more vertex --- nothing is forced.}
\label{fig:floyd}
\end{figure}

\section{State expansion}\label{sec:state}\index{state expansion}
When the question carries a condition --- a coupon, a ticket, a mode
--- the graph is too small: run shortest paths on
\emph{position $\times$ state}. Vertex $(v, \text{used})$ is a
different vertex from $(v, \text{unused})$.
\begin{example}[CSES \emph{Flight Discount}]\label{ex:discount}
One ticket may be halved.
\par\smallskip
\textbf{Strategy.} Double the graph: coupon-unused layer and spent
layer, with a half-price edge between them. Ordinary Dijkstra on the
expanded graph.
\par\smallskip
\begin{solution}
\lstinputlisting{code/flight_discount.cpp}
Real output (CSES sample):
\begin{lstlisting}[numbers=none]
2
\end{lstlisting}
\end{solution}
\end{example}
\begin{figure}[ht]\centering
% fig39-layers.tex -- Figure 7.11: two-layer expansion for one half-price ticket.
% fstyle.tex -- shared TikZ styles for the figures of this book (palette from lifeos.sty).
\tikzset{%
  cell/.style={draw=ink, minimum width=7mm, minimum height=6.5mm, inner sep=0pt, font=\footnotesize, fill=white},
  cellh/.style={cell, fill=accentsoft, draw=accent},
  cella/.style={cell, fill=amberbg, draw=amber},
  cellg/.style={cell, fill=goodbg, draw=good},
  cellx/.style={cell, fill=graybg, draw=slate, text=slate},
  vtx/.style={circle, draw=ink, minimum size=6.5mm, inner sep=0pt, font=\footnotesize, fill=white},
  vtxh/.style={vtx, draw=accent, fill=accentsoft},
  vtxa/.style={vtx, draw=amber, fill=amberbg},
  vtxg/.style={vtx, draw=good, fill=goodbg},
  vtxd/.style={vtx, draw=slate, fill=graybg, text=slate},
  lab/.style={font=\footnotesize, text=slate},
  labA/.style={font=\footnotesize, text=accent},
  labM/.style={font=\footnotesize, text=amber},
  idx/.style={font=\tiny, text=slate},
  arr/.style={-{stealth}, thick, ink},
  arrA/.style={-{stealth}, thick, accent},
  arrM/.style={-{stealth}, thick, amber},
  arrG/.style={-{stealth}, thick, good},
  edge/.style={thick, ink},
  edgeA/.style={thick, accent},
  edgeM/.style={thick, amber},
}%
% \cellrow{prefix}{x0}{y}{style}{v0,v1,...}: a row of cells named prefix0, prefix1, ... (pitch 0.75cm)
\providecommand{\cellrow}[5]{\foreach \v [count=\i from 0] in {#5} {\node[#4] (#1\i) at (#2+\i*0.75,#3) {\v};}}%
% \cellidx{prefix}{n}{start}: index labels under cells prefix0 .. prefix(n-1)
\providecommand{\cellidx}[3]{\foreach \i in {0,...,\numexpr#2-1\relax} {\node[idx, below=0pt of #1\i] {\the\numexpr\i+#3\relax};}}%
% \fitfig{<tikzpicture>}: shrink the picture to \linewidth only when it would overflow.
\ifdefined\fitbox\else\newsavebox{\fitbox}\fi%
\providecommand{\fitfig}[1]{\sbox{\fitbox}{#1}\ifdim\wd\fitbox>\linewidth\resizebox{\linewidth}{!}{\usebox{\fitbox}}\else\usebox{\fitbox}\fi}%

\begin{tikzpicture}[font=\footnotesize]
\node[vtx, minimum size=8mm] (a0) at (0,1.6) {$A$}; \node[vtx, minimum size=8mm] (b0) at (3.2,1.6) {$B$};
\node[vtxd, minimum size=8mm] (a1) at (0,-0.4) {$A'$}; \node[vtxd, minimum size=8mm] (b1) at (3.2,-0.4) {$B'$};
\draw[arr] (a0) -- (b0) node[midway, above, lab] {$w$};
\draw[arr] (a1) -- (b1) node[midway, below, lab] {$w$};
\draw[arrM] (a0) -- (b1) node[pos=0.28, above right=-1pt, labM, fill=white, inner sep=1pt] {$\lfloor w/2\rfloor$};
\node[lab, anchor=west, align=left] at (4.3,1.6) {upper layer: coupon unused};
\node[lab, anchor=west, align=left] at (4.3,-0.4) {lower layer: coupon spent};
\node[lab, anchor=west, text=amber, align=left] at (4.3,0.6) {discount edge: drops a layer only};
\end{tikzpicture}

\caption{The two-layer expansion. Notice the discount edge only
descends --- the coupon can't be unspent.}
\label{fig:layers}
\end{figure}
\begin{example}[USACO-G \emph{Why Did the Cow Cross the Road}]\label{ex:cow}
Crossing into a different label consumes a ticket; each ticket covers
three crossings.
\par\smallskip
\textbf{Strategy.} The budget is a state: $(r,c,k)$ with $k$ the
crossings used on the current ticket. Same-label moves keep $k$;
boundary moves increment it; the third crossing pays the toll.
\par\smallskip
\begin{solution}
\lstinputlisting{code/cowroad.cpp}
Real output (3$\times$3 demo, $T{=}3$):
\begin{lstlisting}[numbers=none]
4
\end{lstlisting}
\end{solution}
\end{example}
A final named variant: \emph{0-1 BFS}\index{0-1 BFS} --- weights in
$\{0,1\}$ only --- replaces the heap with a deque pushing weight-0
edges to the front.

\begin{exercise}[easy]\label{exr:order}
Does relaxation order matter for Dijkstra's correctness? For
Bellman--Ford's round count?
\end{exercise}
\begin{exercise}[medium]\label{exr:earlyexit}
Bellman--Ford may exit early when a round improves nothing. Show the
early exit can never fire while a reachable negative cycle still has
improvable vertices.
\end{exercise}
\begin{exercise}[medium]\label{exr:whypq}
Why a priority queue and not repeated scans for the minimum?
Estimate both costs on $V{=}10^5$, $E{=}2\cdot10^5$.
\end{exercise}
\begin{exercise}[medium]\label{exr:spotstate}
Name the state for: ``shortest route with at most $k$ road changes''
and ``arriving by exactly time $T$ counts differently than late.''
\end{exercise}

\section{Chapter summary}
\begin{table}[ht]\centering
{\small
\begin{tabularx}{\linewidth}{@{}lYYY@{}}
\toprule
\tablehead\thead{Weights} & \thead{Algorithm} & \thead{Cost} & \thead{Also gives}\\
\midrule
unweighted & BFS & $O(V{+}E)$ & layers\\
non-negative & Dijkstra & $O(E\log V)$ & settled order\\
any (incl.\ neg.) & Bellman--Ford & $O(VE)$ & negative cycles\\
all pairs & Floyd--Warshall & $O(V^3)$ & \texttt{next} paths\\
0/1 only & 0-1 BFS & $O(V{+}E)$ & deque order\\
with a condition & expanded graph & $\times$ states & coupon semantics\\
\bottomrule
\end{tabularx}}
\caption{Pick by the weight law and the shape of the question ---
one source, all pairs, or a condition to fold in.}
\label{tab:sp}
\end{table}

\section{Common mistakes}
\index{common mistakes}
\begin{itemize}
\item \textbf{Dijkstra on negative edges.} The settled set is final
--- one negative edge invalidates the permanence and the algorithm
has no repair path.
\item \textbf{Seeding Bellman--Ford from one source when any cycle
counts.} Only reachable-from-$s$ cycles surface; seed 0 everywhere to
catch all.
\item \textbf{Floyd with $k$ inside.} Loop order \emph{is} the
invariant: $k$ outermost, or intermediate vertices may be reused
illegally.
\item \textbf{Deleting queue entries.} Decrease-key is awkward ---
push a better entry and skip stale pops instead.
\item \textbf{Forgetting \texttt{dist[u]}+w can overflow.} Guard
$\texttt{dist[u]} < \texttt{INF}/2$ before adding to avoid wrapping
past unreachable vertices.
\end{itemize}

\section{Additional exercises}
\begin{exercise}[easy] CSES \emph{Shortest Routes I}. \end{exercise}
\begin{exercise}[easy] CSES \emph{Shortest Routes II}. \end{exercise}
\begin{exercise}[easy] USACO-S \emph{Milk Pails} --- BFS on (x,y) fill states. \end{exercise}
\begin{exercise}[medium] CSES \emph{Flight Discount}. \end{exercise}
\begin{exercise}[medium] CSES \emph{Flight Routes} --- $k$ shortest. \end{exercise}
\begin{exercise}[medium] CSES \emph{Investigation} --- count min paths. \end{exercise}
\begin{exercise}[medium] USACO-G \emph{Milk Pumping} --- Dijkstra + max-flow constraint. \end{exercise}
\begin{exercise}[medium] USACO-G \emph{Shortcut} --- tree + shortcut edges. \end{exercise}
\begin{exercise}[hard] CSES \emph{High Score} --- negative cycle on positive edge. \end{exercise}
\begin{exercise}[hard] CSES \emph{Cycle Finding} --- the reconstruction, coded. \end{exercise}
\begin{exercise}[hard] USACO-G \emph{Lasers and Mirrors} --- direction-state BFS. \end{exercise}
\begin{exercise}[hard] USACO-G \emph{A Pie for a Pie} --- two-shelf trade graph. \end{exercise}
\begin{exercise}[hard] USACO-G \emph{Telephone} --- breed-count state Dijkstra. \end{exercise}

\chapternotes
Relaxation-first framing and the cut argument follow
Laaksonen~\cite{laaksonen2024} and Sannemo~\cite{sannemo2025};
state expansion and the coupon/layer trick are standard in the USACO
Guide's Gold modules~\cite{usacoguide} and
cp-algorithms~\cite{cpalgorithms}; 0-1 BFS and the negative-edge
autopsy appear in Halim~\cite{halimcp4}; archive problems are
official CSES/USACO~\cite{cses}.
